\section{Conclusion}

Geo-distributed LLM capacity is request-dependent: a region that is optional
for one request may be indispensable to another.
This makes irreversible arrival-time placement vulnerable to avoidable
allocation-induced SLO misses as later demand changes the opportunity cost of
regional capacity.

Our diagnosis shows that simply delaying commitment is insufficient.
At an 80-ms fixed late-binding delay, removing the direct waiting penalty
recovers only 0.6 percentage points of SLO attainment, whereas making held
requests immediately visible to prospective capacity accounting recovers
15.0 points.
The dominant recoverable loss therefore comes from delaying demand visibility
together with placement finality.

RAVEL separates these lifecycle decisions.
Requests enter capacity planning immediately through tentative placement and
prospective accounting, remain revisable while Router-owned, and become final
only after successful engine handoff.
Replica-specific SLO quotes, bounded pre-handoff replanning, and engine-local
protection make this separation practical.

On our physical-WAN vLLM deployment, RAVEL improves request-level SLO attainment
across conversational, bursty, and multi-stage workloads while retaining
millisecond-scale online planning overhead.
More broadly, preserving online flexibility need not require hiding
already-arrived work from the state used to make future decisions.
